\documentclass{book}
\usepackage[utf8]{inputenc}
\usepackage{amsmath}
\usepackage{amsfonts}
\usepackage{tcolorbox}

\begin{document}


\tableofcontents
\chapter{Recall of Thermodynamics}
\chapter{Thermodynamics of Phase Transitions}
\chapter{Recall of ensembles theory}
\chapter{Statistical mechanics and phase transitions}
This chapter aims to make some general points about how statistical mechanics "frames" phase transitions.
These are very general things that we will encounter throughout and serve as "fixed points" for the reasonings
that we will make onwards. These points are:

\begin{enumerate}
    \item In statistical mechanics, we initially perform calculations with a finite number of particles $N$. However, to recover results consistent with classical thermodynamics, we must take the limit $N \to \infty$. 
    \item Fluctuation-dissipation relation
\end{enumerate}

\section{The bulk free energy: why taking the $N\rightarrow \infty$ limit of stat mech quantities}

\begin{tcolorbox}[colback=blue!5!white,colframe=blue!75!black,title=Key Concept]
In statistical mechanics, we initially perform calculations with a finite number of particles $N$. However, to recover results consistent with classical thermodynamics, we must take the limit $N \to \infty$. 
\end{tcolorbox}

\subsection*{1. The Finite vs. Infinite Gap}
In statistical mechanics, we deal with sums over discrete states. For finite $N$, these sums are always well-behaved and analytic. Thermodynamics, however, describes the \textbf{emergent behavior} of a continuum. Sending $N$ to infinity allows us to observe if these microscopic properties stabilize into the macroscopic values (like pressure or temperature) observed in the real world.

\subsection*{2. Why the Limit Must Be Proved}
The existence of the limit is not guaranteed; it depends entirely on the nature of the particle interactions. The limit might fail to exist if:
\begin{itemize}
    \item \textbf{Long-range interactions:} If forces do not decay quickly enough (e.g., certain gravitational models), the energy could grow faster than the volume, causing the bulk free energy density $f_b$ to diverge.
    \item \textbf{Unstable potentials:} If the short-range attraction is too strong, the system might collapse into a state of infinite density.
\end{itemize}

\subsection*{3. The Stability Condition}
For a statistical mechanics model to be physically valid in the thermodynamic limit, we generally must prove:
\begin{enumerate}
    \item \textbf{Stability:} The potential energy per particle must have a lower bound.
    \item \textbf{Temperateness:} The interaction potential must decay sufficiently fast at large distances.
\end{enumerate}
If these conditions are met, the limit for the bulk free energy density exists:
\[ f_b \equiv \lim_{V \to \infty} \frac{F_\Omega}{V} \]

\subsection*{4. Phase Transitions}
The $N \to \infty$ limit is mathematically essential for describing phase transitions. A finite sum of smooth functions is always smooth; therefore, the sharp "kinks" or singularities associated with phenomena like boiling or freezing only appear mathematically when we reach the infinite limit.




\section{The fluctuation-dissipation relation}
\begin{tcolorbox}[colback=blue!5!white,colframe=blue!75!black,title=Key Concept]
If you want to know how a system reacts to an outside force, just look at how it 'jiggles' on its own.
\end{tcolorbox}
\end{document}